% 01-description.tex - Description and Scope

\chapter{Description}
\label{sec:description}

\section{Purpose}

This document defines the interface control specifications for the \SC{} system. \SC{} is an autonomous data management system that runs on the \MCS{} task processor at LWA stations, providing intelligent file copying and deletion services for the Data Recorders (\DR{}s). Its primary functions include:

\begin{itemize}
    \item Queued file copying from \DR{}s to local and remote destinations, e.g., the \LWA{} \UCF{}
    \item Queued file deletion with configurable purge thresholds
    \item Per-\DR{} queue management with pause and resume capabilities
    \item Automatic retry of failed transfers with configurable backoff
    \item Bandwidth-limited remote transfers
    \item Station monitoring to avoid copying during active observations
\end{itemize}

This ICD serves as the authoritative reference for developers, integrators, and operators working with the \SC{} system.

\section{Scope}

This document covers the following aspects of the \SC{} system:

\begin{enumerate}
    \item \textbf{Command Interface}: UDP-based command protocol for communicating with \SC{}, including system commands, copy/delete commands, and queue management commands.

    \item \textbf{Queue Management}: Per-\DR{} persistent queuing system backed by SQLite, including queue operations, retry logic, and automatic purging.

    \item \textbf{Copy Engine}: The \code{rsync}-based file transfer engine, including local copies, remote copies via SSH, bandwidth limiting, pause/resume, and destination file truncation.

    \item \textbf{Station Monitoring}: Monitoring of the \MCS{} log file to track \DR{} busy/idle states and avoid copying during active recording. Because \SC{} operates independently from \MCS{}, reading the log file is its only mechanism for observing station activity.

    \item \textbf{Status and Error Handling}: System state management, exit codes, error reporting, and email notifications.
\end{enumerate}

\section{System Overview}

\SC{} runs on the \MCS{} task processor at LWA stations and communicates with the \DR{}s via SSH for file operations. Although \SC{} uses the same network infrastructure as \MCS{} and other LWA subsystems, it operates independently: \MCS{} does not see or route \SC{} messages, and \SC{} cannot see \MCS{} or subsystem messages. This separation is why \SC{} must monitor the \MCS{} log file to track \DR{} state (see Section~\ref{sec:station-monitoring}). The system receives commands via UDP and manages independent copy queues for each \DR{}. The system architecture consists of:

\begin{description}
    \item[Command Handler] Receives and processes UDP commands using a packet format based on the \MCS{} packet protocol. The command handler listens for incoming packets, dispatches commands to the appropriate functions, and returns responses.

    \item[Reference Server] A \O{}MQ REP socket that serves unique reference IDs for command packets. The reference ID counter is periodically saved to disk and survives restarts.

    \item[Station Monitor] A background thread that tails the \MCS{} \file{mselog.txt} log file to track \DR{} busy/idle states. Because \SC{} cannot directly observe \MCS{} traffic, the log file is its only source of station activity information. When a \DR{} transitions to idle, its copy queue is automatically resumed; when it becomes busy (recording), the queue is paused.

    \item[DR Queue Managers] One background thread per \DR{} that processes the persistent copy queue. Each manager pulls items from the queue, executes copy or delete operations, handles retries, and manages completed/failed file tracking.

    \item[Copy Engine] The \code{InterruptibleCopy} class that wraps \code{rsync} for file transfers. Supports local copies, remote copies via SSH, bandwidth limiting, pause/resume via process signals, and progress tracking.
\end{description}

\section{Key Functions}

\SC{} performs the following primary functions:

\subsection{Command Processing}

\SC{} accepts commands via UDP packets:
\begin{itemize}
    \item \cmd{PNG}: Ping/heartbeat verification
    \item \cmd{INI}: System initialization
    \item \cmd{SHT}: System shutdown
    \item \cmd{SCP}: Queue a file copy operation
    \item \cmd{SRM}: Queue a file delete operation
    \item \cmd{PAU}: Pause a \DR{} copy queue
    \item \cmd{RES}: Resume a \DR{} copy queue
    \item \cmd{SCN}: Cancel a queued copy
    \item \cmd{RPT}: Report MIB (Management Information Base) entries
\end{itemize}

\subsection{File Transfer}

When a copy command is received, \SC{}:
\begin{itemize}
    \item Assigns a unique ID to the copy operation
    \item Adds the operation to the appropriate \DR{}'s persistent queue
    \item Waits for the \DR{} to become idle before starting the transfer
    \item Executes the transfer using \code{rsync} with progress monitoring
    \item Tracks completed files for later source deletion
    \item Retries failed transfers with configurable wait intervals
\end{itemize}

\subsection{Source File Purging}

\SC{} manages source file cleanup:
\begin{itemize}
    \item Tracks completed transfers and their file sizes
    \item When accumulated completed data exceeds the \code{purge\_size} threshold (default: 1~TB), deletes source files from the \DR{}
    \item Purge runs are performed daily at 18:00~UTC
    \item Spectrometer files are only eligible for deletion after being copied to the remote archive
\end{itemize}

\subsection{Failure Notification}

\SC{} sends email notifications for failed transfers:
\begin{itemize}
    \item Failed file reports are generated daily alongside purge operations
    \item Reports include file paths, sizes, and failure reasons
    \item Notifications are sent to the operations mailing list via SMTP
\end{itemize}

\section{Applicable Stations}

This ICD applies to all LWA stations running the \SC{} system. Station-specific configuration parameters are defined in the \file{defaults.json} configuration file.

